\documentclass[11pt]{article}
\usepackage[utf8]{inputenc}
\usepackage{amsmath,amssymb}
\usepackage[margin=1in]{geometry}
\usepackage{hyperref}
\emergencystretch=3em

\title{The first $\log n$: the two-dimensional standard integral\\
$Z_2(n)\sim\tfrac14 S_{1/2}(a)\,n^{-1/2}\log n$}
\author{Claude, with Daniel Murfet}
\date{2026-09-06}

\begin{document}
\maketitle
\sloppy

\begin{abstract}
The first genuinely multivariate result of the grammar arc and the paper's first $\log n$: for the chart
integral with $d=2$, $k=(1,1)$, $h=(0,0)$, $\xi\equiv a$, $\eta\equiv1$,
\[
Z_2(n)=\int_0^b\!\!\int_0^b e^{-\beta nu^2v^2+\beta\sqrt n\,uv\,a}\,dv\,du
\;\sim\;\frac{S_{1/2}(a)}{4}\,n^{-1/2}\log n\qquad(n\to\infty),
\]
with the explicit bound $|Z_2(n)-\tfrac A2n^{-1/2}\log n|\le(2A|\log b|+E+M)n^{-1/2}$ for $\sqrt n b^2\ge1$,
$A=S_{1/2}(a)/2$. The candidate exponent $\mu=\tfrac12$ appears with multiplicity $m=2$, exactly as the
state density $-(\log\tau)/(4\sqrt\tau)$ of $\tau=u^2v^2$ predicts. No two-dimensional Fubini theorem is
used. Zero \texttt{sorry}/\texttt{axiom}.
\end{abstract}

\section{Setting}
With the one-dimensional theory complete, the next frontier was the $\log n$ powers of
\texttt{thm:TaylorTree}, which arise only for $d\ge2$ from repeated poles of the state density's Mellin
transform. Rather than build the Mellin/state-density machinery, a GPT-6 Astra consult
(\texttt{tide-log/gpt6\_log2d\_v1.md}) suggested a route through two one-dimensional substitutions and a
squeeze: with $f(s)=e^{-\beta s^2+\beta as}$, $F(x)=\int_0^xf$, and the logarithmic primitive
$H(L)=\int_0^LF(x)/x\,dx$, the inner integral is $F(\sqrt n\,ub)/(\sqrt n\,u)$, the outer then gives
$Z_2(n)=n^{-1/2}H(\sqrt n\,b^2)$, and $H(L)=A\log L+O(1)$ uniformly. The logarithm is the integrated
$1/u$ between the saturation scale $u\sim1/(b\sqrt n)$ and the fixed chart endpoint.

\section{The things formalised}
\texttt{QuadraticKernel.lean}: \texttt{quadKernel}, \texttt{quadPrimitive}, \texttt{quadMass}
($=S_{1/2}(a)/2$), \texttt{quadMoment} ($=S_1(a)/2$), \texttt{logPrimitive}; Gaussian domination;
$0\le F(x)\le Ex$; $0\le A-F(x)\le M/x$; $\int_1^Ldx/x^2=1-1/L$; and the squeeze
{\small
\begin{verbatim}
theorem logPrimitive_bounds (β a L : ℝ) (hβ : 0 < β) (hL : 1 ≤ L) :
    quadMass β a * Real.log L - quadMoment β a ≤ logPrimitive β a L ∧
      logPrimitive β a L ≤ quadMass β a * Real.log L + Real.exp (β * a ^ 2 / 2)
\end{verbatim}
}
\texttt{LogTwoD.lean}: \texttt{twoD\_inner}, \texttt{twoDIntegral\_eq} ($Z_2=n^{-1/2}H(\sqrt nb^2)$),
\texttt{twoDIntegral\_sub\_le} (the explicit bound), \texttt{isLittleO\_inv\_sqrt\_log}, and
{\small
\begin{verbatim}
theorem twoDIntegral_isEquivalent (β a b : ℝ) (hβ : 0 < β) (hb : 0 < b) :
    (fun n : ℝ => twoDIntegral β a b n) ~[atTop]
      fun n : ℝ => fluctuation β (1 / 2) a / 4 * (n ^ (-(1 / 2 : ℝ)) * Real.log n)
\end{verbatim}
}

\section{Proof strategy}
Everything is one-dimensional. Mass and moment are the unit-22 kernel identity at $n=1$, $k=1$,
$h=0,1$. The primitive is an interval integral (continuity from
\texttt{intervalIntegral.continuous\_primitive}); $F(x)\le Ex$ by \texttt{integral\_mono\_on} against a
constant; $A-F(x)=\int_x^\infty f\le M/x$ by comparing $xf(s)\le sf(s)$. For $H$, split $(0,L]$ at $1$:
on $(0,1]$ the integrand $F(x)/x\le E$; on $(1,L]$ write $F/x=A/x-(A-F)/x$, integrate $A/x$ exactly to
$A\log L$ (\texttt{integral\_one\_div\_of\_pos}) and bound $(A-F)/x\le M/x^2$ with $\int_1^L dx/x^2\le1$
(\texttt{integral\_zpow}). The two substitutions are \texttt{intervalIntegral.integral\_comp\_mul\_left};
the outer one uses the everywhere-valid identity $F(cu)/u=c\cdot F(cu)/(cu)$ (both sides vanish at $u=0$
under Lean's division convention). Finally $\log(\sqrt nb^2)=\tfrac12\log n+2\log b$, and
$n^{-1/2}=o(n^{-1/2}\log n)$ via $1/\log n\to0$.

\section{Roadblocks and resolutions}
\begin{itemize}
\item \texttt{open intervalIntegral} together with \texttt{MeasureTheory} makes
\texttt{integral\_const\_mul}, \texttt{integral\_sub} ambiguous; keep the \texttt{intervalIntegral.}
prefix explicit.
\item \texttt{integrableOn\_const} now has autoParam side conditions that \texttt{aesop} could not
discharge for \texttt{volume (Ioc 0 L) ≠ ⊤}; pass \texttt{measure\_Ioc\_lt\_top.ne} explicitly.
\item \texttt{intervalIntegrable\_inv}/\texttt{intervalIntegrable\_zpow} are not the installed names;
integrability of $1/x$ and $1/x^2$ on $(1,L]$ was obtained by bounded domination instead (robust, no
name dependence). \texttt{Set.Ioc\_disjoint\_Ioc\_same} likewise absent --- prove disjointness directly.
\item \texttt{simp only [...] at h} on the unit-22 identity mangled the integrand before it could be
matched; rewrite the integrand to the kernel first (\texttt{setIntegral\_congr\_fun}), then simplify the
constants with \texttt{norm\_num}.
\item Pointwise goals of \texttt{IntegrableOn.congr\_fun} arrive un-beta-reduced: \texttt{beta\_reduce}
before \texttt{rw}. \texttt{Integrable.const\_mul}/\texttt{.sub} return a bare \texttt{Integrable}
(an \texttt{And}) --- type-ascribe as \texttt{IntegrableOn} before \texttt{.congr\_fun}/\texttt{.mono\_set}.
\item \texttt{rw [...] at *} silently left the squeeze hypothesis in its old form; rewrite the log
identity into that hypothesis explicitly so that \texttt{linarith} sees identical atoms.
\item In the last step, \texttt{← Real.rpow\_neg} found no \texttt{(n\^{}y)⁻¹} because the goal had a
division; rewrite \texttt{Real.rpow\_neg} forwards on the other side and let \texttt{ring} finish.
\end{itemize}

\section{What was Mathlib, what was new}
Mathlib: \texttt{intervalIntegral.integral\_comp\_mul\_left}, \texttt{integral\_of\_le},
\texttt{continuous\_primitive}, \texttt{integral\_mono\_on}, \texttt{integral\_one\_div\_of\_pos},
\texttt{integral\_zpow}, \texttt{Real.log\_sqrt}, \texttt{tendsto\_log\_atTop},
\texttt{Tendsto.inv\_tendsto\_atTop}, \texttt{IsEquivalent.congr\_right}. New: the whole
logarithmic-primitive theory and the first formal $\log n$ asymptotic for a standard integral. The
leading constant was checked numerically two ways ($S_{1/2}(a)/4=0.512920$ against the erf closed form)
and the squeeze verified at $n=10^2,\dots,10^8$ before any Lean was written.

\section{Lessons}
When an expansion is expected to carry a logarithm, look for a one-dimensional reduction in which the
logarithm is literally $\int dx/x$ over a growing range; the uniform $O(1)$ squeeze of the logarithmic
primitive is much cheaper than any meromorphic-continuation argument, and it needs no Fubini theorem
because the inner integral is a primitive.

\section{Follow-ups}
The exact next-order term ($-B\,n^{-1/2}$ with the log-moment $B=\int f\log s$), the product-measure
form of $Z_2$ via Fubini, general $(k_1,k_2)$ and $(h_1,h_2)$ in two dimensions (multiplicity $2$ iff the
two candidate exponents coincide), and the general-$d$ multiplicity structure.
\end{document}
